\begin{proposition}\label{prop:nonic-no230}
No elliptic curve over a number field of degree $9$ has a subgroup
isomorphic to $C_2\times C_{30}$ or $C_{60}$.
\end{proposition}
\begin{proof}
Put $X=X_1(2,30)$ and $E_0=X_1(15)$. On the noncuspidal locus,
$X$ classifies an elliptic curve with a point of order $15$ and an
ordering of its three nonzero points of order $2$. Quotienting by the
alternating subgroup gives maps
\[
 \pi:X\longrightarrow Y=X/A_3,\qquad f:Y\longrightarrow E_0.
\]
For the universal Tate curve $\mathcal E_{15}$, the function field is
\[
 \mathbb Q(Y)=\mathbb Q(E_0)(\rho),\qquad
 \rho^2=\Delta(\mathcal E_{15}).
\]
Indeed, after completing the square, the expression
$4(e_1-e_2)(e_2-e_3)(e_3-e_1)$ in the three roots of the
$2$-division cubic has stabilizer $A_3$ and square equal to the
elliptic discriminant. All sixteen geometric cusps of $E_0$ have odd
width. Thus $f$ ramifies at every cusp, and $Y$ has genus $9$.
All these models and maps have good reduction at $7$.

The rational Galois orbits of cusps on $Y$, in the ordering of the
certificate, are $C_1,\ldots,C_8$ with degrees
\[
 (1,1,4,2,4,1,2,1).
\]
Each is the unique lift of the corresponding cusp orbit on $E_0$.
We identify these orbits in characteristic zero and at $7$ using the
valuation pairs of the modular units $b^4/\Delta$ and $b$ on the
universal Tate curve. There are eight distinct pairs on $E_0$ and
sixteen on $X_1(30)$, equal to their theoretical numbers of rational
Galois orbits. These units have no noncuspidal or vertical divisor
away from the level. Hence the identifications commute with reduction.

By \cite[Theorem~3.1, Proposition~4.16 and Table~3]{DEHMZ},
$J_X(\mathbb Q)$ has rank zero and is generated by rational cuspidal
divisors. Its invariants are
$[2,2,2,2,2,24,8160,8160]$. The rank of $J_Y(\mathbb Q)$ is also
zero, since $\pi_*\pi^*=[3]$.
Choose a rational cusp $c_0$ on $Y$, and put
\[
 H=\langle[C_i-(\deg C_i)c_0]:1\leq i\leq8\rangle.
\]
For any rational effective divisor $D_X$ of degree $9$ on $X$, the
class of $\pi_*D_X-9c_0$ belongs to $H$. This uses cuspidal generation
on the source; no cuspidal generation assumption on $J_Y(\mathbb Q)$
is needed. A proved finite class-group computation gives
\[
 J_Y(\mathbb F_7)=[2,2,2,2,2,2,102,10608],\qquad
 H=[2,2,2,2,2,8],\qquad \#H=256.
\]
Odd unramified torsion injectivity identifies $H$ with its reduction.

We enumerate all $956$ rational effective cusp divisors of degree $9$
on $Y$. Their normalized classes exhaust the $256$ elements of $H$.
For each class choose one such divisor $L=\sum_i v_iC_i$. Write
\[
 L=f^*B+C,\qquad
 B=\sum_i\lfloor v_i/2\rfloor f(C_i),\qquad
 C=\sum_i(v_i\bmod2)C_i.
\]
Pullback gives at least $\deg B$ sections if $\deg B>0$, and the
constant section if $B=0$. Exact Riemann--Roch calculations over
$\mathbb F_7$ give the following exhaustive alternatives:
\[
\begin{array}{c|r}
\text{alternative}&\text{number of classes}\\\hline
h^0(L)=1,\text{ with an effective cuspidal representative}&84\\
H^0(Y,L)=H^0(E_0,B),\quad C>0&144\\
\text{remaining pencils, }h^0_{\mathbb F_7}(L)=2&28
\end{array}
\]
In the second row the pullback lower bound equals the finite
Riemann--Roch upper bound. The equality therefore holds in
characteristic zero as well. Since $\deg L=9$ is odd, every such
pullback series has a positive fixed cuspidal part $C$.
The first two rows cannot contain a divisor supported entirely on
the noncuspidal locus.

Suppose that $E/K$ has the proposed subgroup, with $[K:\mathbb Q]=9$.
Its modular point, traced with multiplicity if necessary, gives a
noncuspidal rational effective divisor $D_X$ of degree $9$ on $X$.
Entirely cuspidal reduction at $7$ is impossible. Indeed, lift such a
reduction to a cusp divisor $C_X$ over an unramified cyclotomic
extension of $\mathbb Q_7$. The class $[D_X-C_X]$ is torsion by the
rank-zero assertion and Manin--Drinfeld. It reduces to zero and thus
is zero by odd unramified torsion injectivity. Since $D_X\ne C_X$,
this implies $h^0(D_X)\geq2$ after extension, and hence over
$\mathbb Q$. This contradicts
$\operatorname{gon}_{\mathbb F_7}X\geq10>9$ from
\cite[Table~1]{NV} and specialization of gonality.

It remains to consider the $28$ pencils. We enumerate all eight
projective combinations of an exact finite Riemann--Roch basis for
each pencil, giving $224$ members. At a noncuspidal place on $Y$, the
$2$-division cubic is either split or irreducible, since its
discriminant is a square. In the split case, a support multiplicity
$m$ is distributed in every possible way among the three oriented
source lifts. In the irreducible case it must be divisible by $3$,
and the lift has residue degree multiplied by $3$.
The orientation is selected by
\[
 \rho=4(e_1-e_2)(e_2-e_3)(e_3-e_1)/u^6,
\]
where $u$ is the scaling to the universal Tate normal form.
Every source level, residue degree and normalization is checked
exactly. Cuspidal source complements are left unrestricted, which
only enlarges the necessary-condition set. Exactly $42$
noncuspidal source reductions remain: $24$ single places of degree
$9$ and $18$ single places of degree $8$ with a cuspidal complement
of degree $1$.

Apply the three rational maps from $X$ to $X_1(30)$ retaining,
respectively,
\[
 Q_{30},\qquad P_2+16Q_{30},\qquad P_2+Q_{30}.
\]
The three points have order $30$ because $P_2\ne15Q_{30}$.
By \cite[Theorem~4.13 and Table~2]{DEHMZ}, the rational Jacobian of
$X_1(30)$ is cuspidal with invariants $[4,8160]$. Its finite
reduction at $7$ has invariants $[4,4,4,136,16320]$, and the rational
cuspidal image has index $4352$.
For a source reduction of degree $9$, the normalized target class
must belong to that image. For a source reduction of degree $8$,
we explicitly allow every degree-one finite cusp as a complement.
Their classes give nine distinct shifts modulo the rational image.
All norm multiplicities are retained in the pushforward.

The $126$ actual maps give $21$ distinct target classes. Every one
of the $42$ source reductions has a target class failing its
necessary image condition. Independent integer Hermite-normal-form
arithmetic checks every subgroup and shifted-coset test.
A separate finite-field implementation checks all $78$ intermediate
primitive source levels, all $126$ level maps, and the complete
$224$-member lift enumeration. This excludes $D_X$ and proves the
assertion for $C_2\times C_{30}$.

Finally, a point $P$ of order $60$ gives, after quotienting by
$\langle30P\rangle$, a point of order $30$ and an independent rational
point of order $2$ in the dual kernel. Their independence follows
because the dual isogeny sends the order-$2$ multiple of the first
point to $30P\ne0$, while it kills the dual kernel. This would give
$C_2\times C_{30}$, already excluded.
\end{proof}
